% =============================================================================
% Chapter 1: Introduction
% =============================================================================

\chapter{Introduction}
\label{ch:introduction}

Cloud computing has fundamentally changed how organizations deploy and manage software. Containers, orchestrated by platforms such as Kubernetes, have become the standard unit of deployment across industries ranging from financial services to scientific computing. As organizations mature in their cloud adoption, they increasingly operate not one but \emph{multiple} Kubernetes clusters---spread across data centers, cloud providers, and edge locations---to achieve geographic proximity to users, regulatory compliance, fault isolation, and team autonomy.

However, this multi-cluster paradigm introduces a persistent inefficiency: \textbf{resource fragmentation}. At any given moment, some clusters sit idle with surplus CPU and memory, while others struggle under peak demand. Traditional approaches treat each cluster as an isolated silo, forcing operators to over-provision capacity in every cluster to handle worst-case scenarios. Studies on cluster utilization in large-scale deployments consistently show that average resource utilization remains between 20\% and 50\%~\cite{google-borg, azure-traces}, representing a significant waste of computing infrastructure and energy.

The vision of a \textit{computing continuum}---where resources flow seamlessly from edge devices through on-premise clusters to public clouds---has gained traction in both industry and the European research community, notably through initiatives such as the FLUIDOS project~\cite{fluidos}. Yet achieving this vision requires mechanisms that can \emph{discover}, \emph{match}, and \emph{allocate} resources across cluster boundaries in a secure and consistent manner.

This thesis presents a system that addresses this gap: a \textbf{multi-cluster Kubernetes resource sharing platform} with a centralized broker, lightweight per-cluster agents, mutual TLS security, and automatic cluster peering through Liqo.


% =============================================================================
\section{Problem Statement}
\label{sec:problem-statement}

Despite the proliferation of multi-cluster Kubernetes deployments, there is currently no standardized, lightweight mechanism for sharing resources \emph{between} clusters. The key challenges are:

\begin{enumerate}
    \item \textbf{Lack of global visibility.} Each cluster's scheduler sees only its own nodes. There is no federation-wide view of which clusters have spare capacity and which are under pressure.

    \item \textbf{No reservation mechanism.} Even when spare capacity is identified, there is no protocol to \emph{reserve} it for a remote requester while preventing other consumers from claiming the same resources---a classic double-booking problem.

    \item \textbf{Security at cluster boundaries.} Clusters are typically managed by different teams or organizations. Any inter-cluster communication must be mutually authenticated and authorized, without requiring clusters to share credentials or join a common identity provider.

    \item \textbf{Bridging decision and action.} Deciding \emph{which} cluster should provide resources is only half the problem. The system must also establish actual network connectivity between the requester and provider so that workloads can be offloaded.
\end{enumerate}

Existing solutions address parts of this problem. Kubernetes Federation (KubeFed) enables multi-cluster workload distribution but does not provide a resource marketplace or reservation semantics. Liqo enables peer-to-peer cluster peering and pod offloading but relies on manual or static peering decisions rather than dynamic resource matching. Neither system prevents race conditions when multiple requests target the same provider simultaneously.


% =============================================================================
\section{Objectives}
\label{sec:objectives}

The objective of this thesis is to design, implement, and evaluate a system that enables automatic, secure, and consistent resource sharing across multiple Kubernetes clusters. Specifically, the system must:

\begin{enumerate}
    \item \textbf{Collect resource advertisements.} Design a centralized broker that receives and stores real-time resource metrics (CPU, memory, GPU) from all participating clusters.

    \item \textbf{Match reservations to providers.} Implement a decision engine that evaluates available clusters based on a scoring algorithm and selects the optimal provider for each resource request.

    \item \textbf{Prevent double-booking.} Introduce a resource locking mechanism (the \textit{Reserved field}) that ensures consistency under concurrent reservation requests.

    \item \textbf{Secure all communication.} Use mutual TLS (mTLS) with certificate-based identity, where each cluster's identity is derived from its certificate Common Name (CN), eliminating the need for separate authentication tokens.

    \item \textbf{Deploy lightweight agents.} Build a per-cluster agent that monitors local resources, publishes advertisements, receives instructions, and consumes minimal CPU and memory.

    \item \textbf{Automate cluster interconnection.} Integrate with Liqo to automatically establish network peering between requester and provider clusters after the broker makes a decision, creating virtual nodes that enable transparent workload offloading.

    \item \textbf{Evaluate the system.} Conduct experiments measuring scalability, latency, concurrency handling, decision accuracy, and resource tracking correctness.
\end{enumerate}


% =============================================================================
\section{Contributions}
\label{sec:contributions}

This thesis makes the following contributions:

\begin{itemize}
    \item \textbf{Centralized Broker Architecture.} A Kubernetes-native broker that collects cluster advertisements, processes reservation requests through a scoring-based decision engine, and manages the complete reservation lifecycle (Pending $\to$ Reserved $\to$ Active $\to$ Released). The broker is implemented as a controller-runtime operator with a REST API secured by mTLS.

    \item \textbf{Reserved Field Mechanism.} A race-condition prevention mechanism where the broker atomically updates a \texttt{Reserved} field in the provider's \texttt{ClusterAdvertisement} before confirming a reservation. This ensures that concurrent requests see reduced availability and cannot double-book the same resources. The agent preserves this field during advertisement updates through a two-step read-then-write protocol.

    \item \textbf{Scoring-Based Decision Engine.} A cluster selection algorithm that filters candidates by resource sufficiency and active status, then scores them based on projected post-reservation utilization. The algorithm favors clusters that retain the most headroom after fulfilling the request, distributing load across the federation.

    \item \textbf{mTLS Identity Model.} A security model where cluster identity is derived directly from the client certificate Common Name, managed by cert-manager. This eliminates separate authentication infrastructure (tokens, API keys, OAuth) and provides mutual authentication with minimal operational overhead.

    \item \textbf{Automatic Liqo Peering.} An integration where the agent automatically triggers \texttt{liqoctl peer} upon receiving a \texttt{ReservationInstruction}, establishing encrypted WireGuard tunnels and virtual nodes without manual intervention.

    \item \textbf{Comprehensive Evaluation.} Seven automated test scenarios measuring broker scalability (up to 100 clusters), reservation latency (sub-second end-to-end with synchronous flow), concurrent request handling (zero double-bookings), decision accuracy (83--100\%), resource exhaustion tracking, and agent startup time breakdown.

    \item \textbf{Open-Source Implementation.} The complete system is available as open-source software, implemented in Go with Kubernetes CRDs, suitable for both research and practical deployment.
\end{itemize}


% =============================================================================
\section{Thesis Structure}
\label{sec:thesis-structure}

The remainder of this thesis is organized as follows:

\begin{itemize}
    \item \textbf{Chapter~\ref{ch:background}: Background and Related Work.} Provides the technical background on Kubernetes resource management, Custom Resource Definitions, the controller-runtime framework, multi-cluster architectures, and existing solutions including KubeFed, Liqo, and the FLUIDOS project. Identifies the limitations of current approaches that motivate this work.

    \item \textbf{Chapter~\ref{ch:design}: System Design.} Defines the functional and non-functional requirements, presents the hub-and-spoke architecture, describes the CRD-based data model, the resource calculation formula, the decision algorithm, the REST communication protocol, and the mTLS security model.

    \item \textbf{Chapter~\ref{ch:implementation}: Implementation.} Details the technology stack, project structure, and implementation of the broker (API server, authentication middleware, reservation controller, decision engine) and agent (resource collector, HTTP transport, synchronous reservation flow, instruction polling, Liqo integration).

    \item \textbf{Chapter~\ref{ch:evaluation}: Evaluation.} Presents seven experimental tests conducted on a 96-core server, measuring scalability (with resource exhaustion verification), bandwidth, agent footprint, reservation latency, concurrent reservation handling, decision accuracy, and agent startup time breakdown.

    \item \textbf{Chapter~\ref{ch:conclusions}: Conclusions.} Summarizes the contributions, discusses limitations, and outlines directions for future work.
\end{itemize}
